%%%%%%%%%%%%%%%%%%%%%%%%%%%%%%%%%%%%%%%%%%%%%%%%%%%%%%%%%%%%%%%%%%%%%%%%
% Plantilla TFG/TFM
% Universidad de Murcia. Facultad de Informática
%%%%%%%%%%%%%%%%%%%%%%%%%%%%%%%%%%%%%%%%%%%%%%%%%%%%%%%%%%%%%%%%%%%%%%%%

\chapter{Análisis de objetivos y metodología}
\label{cap:analisis}
Este capítulo presenta los objetivos que orientan el trabajo y la metodología seguida para diseñar, implementar y validar el sistema de verificación de noticias basado en arquitecturas RAG.

\section{Objetivos Generales}

La proliferación de desinformación en el ámbito digital exige herramientas de verificación rápidas y precisas. Los LLMs son herramientas potentes para el procesamiento del lenguaje, pero su tendencia a la ``alucinación'' (la generación de datos plausibles pero falsos) limita su uso en tareas de \textit{Fact-Checking}. Este Trabajo Fin de Grado se centra en el diseño y evaluación de un sistema de Generación Aumentada por Recuperación (RAG) que proporciona al modelo de lenguaje evidencias extraídas de un corpus de noticias veraces, minimizando así la invención de datos.

Para alcanzar esta meta, el proyecto se desglosa en cinco objetivos específicos:

\begin{itemize}
    \item \textbf{Implementar una infraestructura de recuperación robusta:} Desplegar un motor de búsqueda basado en Elasticsearch para indexar, gestionar y consultar de forma eficiente todas los documentos válidos de un corpus de aproximadamente 800.000 noticias.
    \item \textbf{Comparar estrategias de búsqueda léxica y semántica:} Evaluar el rendimiento del algoritmo tradicional BM25 frente a representaciones vectoriales densas mediante el modelo BGE-M3, identificando cuál responde correctamente a más preguntas sobre el conjunto de noticias.
    \item \textbf{Integrar modelos de lenguaje de distinta escala:} Comparar el desempeño de un modelo local pequeño (Small Language Model - Phi-3) frente a LLMs accesibles vía API (Llama 3.3 70B y Qwen 3 32B), comparando la cantidad de respuestas que responden correctamente y analizando ejemplos puntuales de preguntas de diferentes tipos.
    \item \textbf{Establecer un marco de evaluación automatizado:} Diseñar una metodología de evaluación usando ``LLM-as-a-judge'' con un modelo de parámetros masivos (120B) que permita automatizar, en cierta medida, la validación de respuestas sin intervención humana constante.
    \item \textbf{Validar el sistema a gran escala:} Someter la arquitectura final (sistema de recuperacion con mejor desempeño \& modelo de lenguaje con mejor desempeño) a una batería de 220 preguntas complejas de temas variados (\textit{Ground Truth}) clasificadas en 3 tipos (factuales, inferencia y trampa) para obtener métricas de rendimiento con respecto a los otros casos y analizar mejorías.
\end{itemize}

De estos objetivos se desprende la cuestión vertebral de la investigación:

\begin{quote}
    \textit{¿En qué medida la implementación de una arquitectura RAG permite mitigar las alucinaciones en los modelos de lenguaje, y cómo influye la selección del motor de recuperación y la escala del generador en la veracidad de las respuestas obtenidas?}
\end{quote}

\section{Metodología}

La metodología del proyecto se consolidó finalmente en siete fases secuenciales, abarcando todo el ciclo de vida desde el procesamiento de los datos hasta la obtención de las métricas. A continuación se describe este flujo de trabajo:

\begin{enumerate}
    \item \textbf{Ingesta y normalización del corpus (ETL):} Extracción y limpieza de los datos de noticias para su indexación en Elasticsearch. En esta fase se define la estructura de los documentos para facilitar recuperaciones posteriores.
    \item \textbf{Desarrollo del motor de recuperación (Retrieval):} Implementación paralela de la búsqueda por palabras clave (BM25) y la búsqueda vectorial (kNN), validando la relevancia de los documentos recuperados frente a consultas de prueba.
    \item \textbf{Despliegue e integración del flujo RAG:} Configuración del modelo local (Phi-3) y desarrollo de la lógica de conexión entre el recuperador y el generador. En esta etapa se establecen los protocolos de comunicación y los \textit{system prompts} para restringir la respuesta al contexto.
    \item \textbf{Configuración del generador avanzado y diseño del marco de evaluación automatizada (LLM-as-a-Judge):} Integración del modelo \textit{Llama 3.3 70B} mediante API y despliegue de un tribunal basado en un modelo de parámetros masivos (120B). Se ejecuta una prueba piloto de 80 preguntas para validar la consistencia del criterio de evaluación binaria y realizar la primera comparativa entre el modelo local y el de gran escala.
    \item \textbf{Ampliación del marco experimental y selección del generador:} Incorporación de un segundo modelo de lenguaje avanzado (Qwen 3 32B) para evaluar su desempeño frente a Llama 3.3 70B y Phi-3. Esta etapa tiene como fin disponer de una comparativa entre diversas arquitecturas, permitiendo identificar el motor de razonamiento óptimo para su integración en la arquitectura RAG definitiva.
    \item \textbf{Ejecución masiva y control de errores:} Procesamiento del \textit{dataset} completo de 220 preguntas. Se implementan puntos de control debido al elevado tiempo de ejecución y la limitación de recursos computacionales en ciertos momentos.
    \item \textbf{Análisis comparativo y síntesis del sistema óptimo:} Consolidación de resultados en una visualización del conjunto de casos, identificando los componentes ganadores para el ensamblaje y validación de la arquitectura RAG final.
\end{enumerate}